\section{RNG Economy}\label{sec:rng_economics}

\subsection{Token Utility}\label{sec:token_utility}
The \$RNG token serves as the native utility token of the RandAO ecosystem, designed to capture the value of decentralized random number generation. This utility is enforced through protocol design: all randomness requests must be paid for in \$RNG tokens, establishing a foundational and continuous demand for the token.

\subsection{Supply and Demand}\label{sec:supply_demand}
The RNG economy is driven by two primary mechanisms:

\begin{itemize}
    \item \textbf{Demand Side}: Users requiring verifiable randomness must acquire and spend \$RNG tokens to access the service. This creates natural market demand that scales with the protocol's adoption across various use cases such as gaming, NFTs, and blockchain consensus.
    
    \item \textbf{Supply Side}: \$RNG tokens flow to providers through two channels:
    \begin{itemize}
        \item Direct service payments from users requesting randomness
        \item Staking rewards for maintaining collateral in the network
    \end{itemize}
\end{itemize}

\subsection{Market Dynamics}\label{sec:market_dynamics}
When randomness is requested, users can select any subset of providers within the open network. This flexibility creates a competitive marketplace where:

\begin{itemize}
    \item Providers compete on service quality and reliability
    \item Market forces determine efficient pricing for randomness generation
    \item \$RNG tokens are distributed proportionally to participating providers
\end{itemize}

\subsection{Staking Mechanism}\label{sec:staking}
The protocol implements a robust staking mechanism:

\begin{itemize}
    \item \textbf{Collateral Requirement}: Providers must stake wAR to participate
    \item \textbf{Security}: Staked assets ensure economic security against malicious behavior
    \item \textbf{Rewards}: Providers earn \$RNG tokens as staking rewards
    \item \textbf{Alignment}: The dual reward structure (service fees + staking rewards) creates long-term alignment between providers and the protocol
\end{itemize}
